\hypertarget{appendix-f52f8c5a3afd}{}
\section*{Internal and tail padding}
\section{Internal and tail padding}\label{internal-and-tail-padding}

\textbf{Curriculum depth:} Introduced · \textbf{First seen:} C003 · \textbf{Latest development:} C003

\subsection{Definition and mechanism}\label{definition-and-mechanism}

Internal padding fills holes before members to satisfy alignment. Tail padding follows the final member so array stride preserves complete-object alignment. Under 1/4/1 rules, byte--uint32--byte occupies offsets 0/4/8 with end nine and size 12. Padding is not a semantic format field and cannot be relied on as stable encoded content.

\subsection{Prerequisites}\label{prerequisites}

\begin{itemize}
\tightlist
\item
  \hyperref[appendix-a1f7c53b4742]{Alignment and padding} --- Address divisibility determines legal placement and layout holes.
\end{itemize}

\subsection{Used by}\label{used-by}

\begin{itemize}
\tightlist
\item
  \href{}{Field reordering and AoS versus SoA (05-Engine-I, appendix-5522ad176ca0)} --- uses this mechanism as stated in its prerequisites.
\item
  \hyperref[appendix-0682dbfc5f0e]{Packed structures and misaligned access} --- uses this mechanism as stated in its prerequisites.
\end{itemize}

\subsection{Recall and next depth}\label{recall-and-next-depth}

Explain why stride nine breaks the second element's integer alignment.

\textbf{Pending:} Special member overlap, inheritance, and padding-safe comparisons.

\subsection{Source and navigation}\label{source-and-navigation}

\href{Continuous\%20Notes/Computing\%20Foundations\%20-\%20From\%20Bits\%20to\%20C++\%20and\%20Systems\%20APIs\#Where\%20the\%20bytes\%20live}{C003 teaching chapter} · \hyperref[appendix-fd090d982a15]{Part definitions} · \href{}{Dependency Map (02-Foundations-II, appendix-4ad47bdd269a)}
